% ==============================================================================
\chapter{Theory: Parameterized Continuum Mechanics}
\label{ch:theory}
% ==============================================================================

\begin{tcolorbox}[colback=lightbg, colframe=primaryblue, title=\textbf{Chapter at a Glance: Continuum Mechanics \& Shape Parameterization}, fonttitle=\bfseries]
    \begin{description}[leftmargin=!, labelwidth=3.4cm]
        \item[\textbf{Physical Domain}] Parameterized configuration $\Omega(\boldsymbol{\mu}) \subset \mathbb{R}^2$ with design parameters $\boldsymbol{\mu} = (x_c, r)^T$ (notch location, radius).
        \item[\textbf{Governing Law}] Linear elastodynamics balance of momentum $\rho \ddot{\vect{u}} - \nabla \cdot \boldsymbol{\sigma} = \vect{f}_b$ subject to Dirichlet clamping and Neumann tractions.
        \item[\textbf{Core Bottleneck}] Domain variation $\Omega(\boldsymbol{\mu})$ destroys vector space compatibility across parameter instances, preventing direct snapshot POD extraction.
    \end{description}
\end{tcolorbox}

Developing reduced order models for parameterized structures requires a firm foundation in continuum elastodynamics and formal shape parameterization. This chapter establishes the strong and weak variational formulations of linear elasticity, defines the parameter space $\boldsymbol{\mu} = (x_c, r)^T$, and computes the mathematical challenges introduced by domain variations.

% ------------------------------------------------------------------------------
\section{Geometry Configurations \& Shape Parameterization}
% ------------------------------------------------------------------------------

This section defines the parameterized domain geometry $\Omega(\boldsymbol{\mu})$ and establishes how shape variations affect the physical configuration of the notched beam.

Let $\Omega(\boldsymbol{\mu}) \subset \mathbb{R}^2$ define the parameterized physical configuration of a structural continuum. As illustrated in Figure \ref{fig:notched_beam_geometry}, we consider a 2D cantilever beam featuring symmetric circular notches along its upper and lower boundaries.

The shape variation is controlled by a parameter vector $\boldsymbol{\mu} = (x_c, r)^T \in \mathcal{D} \subset \mathbb{R}^2$:
\begin{itemize}[leftmargin=*]
    \item \textbf{Notch Position ($x_c$)}: The longitudinal starting distance of the notch center from the clamped left support ($0.5\,\text{m} \le x_c \le 4.5\,\text{m}$).
    \item \textbf{Notch Radius ($r$)}: The curvature radius of the symmetric boundary cutouts ($0.1\,\text{m} \le r \le 0.45\,\text{m}$).
\end{itemize}

\begin{figure}[htbp]
    \centering
    \begin{tikzpicture}[scale=1.0, >=stealth]
        % Styles
        \tikzset{
            deformed/.style={thick, draw=primaryblue, fill=lightbg, fill opacity=0.8},
            dimension/.style={primaryblue, <->, >=latex, line width=0.8pt},
            dimline/.style={primaryblue, dashed, line width=0.5pt}
        }

        % Left Wall Hatching (Clamped Support)
        \draw[line width=1.5pt] (0,0.8) -- (0,-0.8);
        \foreach \y in {-0.7,-0.5,...,0.7} {
                \draw (0,\y) -- (-0.2,\y-0.15);
            }
        \node[left] at (-0.25, 0) {$\Gamma_D(\boldsymbol{\mu})$};

        % Parameterized Configuration \Omega(\boldsymbol{\mu})
        \draw[deformed] (0,0.6) -- (3.0,0.6) .. controls (3.4,0.6) and (3.6,0.2) .. (3.8,0.2) .. controls (4.0,0.2) and (4.2,0.6) .. (4.6,0.6) -- (6,0.6) -- (6,-0.6) -- (4.6,-0.6) .. controls (4.2,-0.6) and (4.0,-0.2) .. (3.8,-0.2) .. controls (3.6,-0.2) and (3.4,-0.6) .. (3.0,-0.6) -- (0,-0.6) -- cycle;
        \node[primaryblue] at (1.2, 0.0) {\textbf{$\Omega(\boldsymbol{\mu})$}};

        % Dimension: x_c
        \draw[dimline] (0, 0.6) -- (0, 1.1);
        \draw[dimline] (3.8, 0.2) -- (3.8, 1.1);
        \draw[dimension] (0, 1.0) -- (3.8, 1.0) node[midway, above] {$x_c$};

        % Dimension: r
        \draw[dimline] (3.0, 0.6) -- (3.8, 0.6);
        \draw[primaryblue, ->, >=latex, line width=0.8pt] (3.8, 0.6) -- (3.8, 0.2) node[midway, right] {$r$};
    \end{tikzpicture}
    \caption{Schematic of the parameterized notched cantilever beam configuration $\Omega(\boldsymbol{\mu})$ showing Dirichlet boundary support $\Gamma_D(\boldsymbol{\mu})$, notch location $x_c$, and radius $r$.}
    \label{fig:notched_beam_geometry}
\end{figure}

The structural parameterization illustrated above (Figure \ref{fig:notched_beam_geometry}) defines the design space $\mathcal{D}$ via notch position $x_c$ and cutout radius $r$.

\newpage
% ------------------------------------------------------------------------------
\section{Governing Elastodynamic Equations \& Kinematics}
% ------------------------------------------------------------------------------

This section presents the strong form balance of linear momentum and Cauchy stress-strain relationships under small-strain elastodynamics.

Kinematic deformation from the initial configuration $\Omega(\boldsymbol{\mu})$ to the current state under external body forces $\vect{f}_b$ and boundary tractions $\vect{h}_N$ is illustrated below (Figure \ref{fig:continuum_kinematics}).

\begin{figure}[htbp]
    \centering
    \makebox[\textwidth][c]{%
        \begin{tikzpicture}[scale=0.8, >=stealth]
            % Styles
            \tikzset{
                reference/.style={thick, draw=black!70, fill=gray!10},
                deformed/.style={thick, draw=primaryblue, fill=lightbg, fill opacity=0.7},
                dimension/.style={primaryblue, <->, >=latex, line width=0.8pt},
                dimline/.style={primaryblue, dashed, line width=0.5pt},
                force/.style={->, line width=1.5pt, red!70!black},
                displ/.style={->, line width=1.2pt, green!60!black}
            }

            % --- PANEL 1: UNDEFORMED CONFIGURATION ---
            \begin{scope}[xshift=0cm]
                \draw[line width=1.5pt] (0,0.8) -- (0,-0.8);
                \foreach \y in {-0.7,-0.5,...,0.7} {
                        \draw (0,\y) -- (-0.2,\y-0.15);
                    }
                \node[left] at (-0.25, 0) {$\Gamma_D(\boldsymbol{\mu})$};

                \draw[reference] (0,0.6) -- (2.2,0.6) .. controls (2.7,0.6) and (2.9,0.25) .. (3.25,0.25) .. controls (3.6,0.25) and (3.8,0.6) .. (4.3,0.6) -- (6,0.6) -- (6,-0.6) -- (4.3,-0.6) .. controls (3.8,-0.6) and (3.6,-0.25) .. (3.25,-0.25) .. controls (2.9,-0.25) and (2.7,-0.6) .. (2.2,-0.6) -- (0,-0.6) -- cycle;
                \node[black] at (1.2, 0.0) {$\Omega(\boldsymbol{\mu})$};

                \coordinate (X) at (4, 0.0);
                \fill[black] (X) circle (1.5pt) node[above] {$\vect{x}$};
                \node[below, font=\bfseries] at (2.4, -1.8) {(a) Undeformed Configuration};
            \end{scope}

            % --- PANEL 2: DEFORMED CONFIGURATION ---
            \begin{scope}[xshift=8.5cm]
                \draw[line width=1.5pt] (0,0.8) -- (0,-0.8);
                \foreach \y in {-0.7,-0.5,...,0.7} {
                        \draw (0,\y) -- (-0.2,\y-0.15);
                    }
                \node[left] at (-0.25, 0) {$\Gamma_D(\boldsymbol{\mu})$};

                \draw[dashed, gray!40] (0,0.6) -- (2.2,0.6) .. controls (2.7,0.6) and (2.9,0.25) .. (3.25,0.25) .. controls (3.6,0.25) and (3.8,0.6) .. (4.3,0.6) -- (6,0.6) -- (6,-0.6) -- (4.3,-0.6) .. controls (3.8,-0.6) and (3.6,-0.25) .. (3.25,-0.25) .. controls (2.9,-0.25) and (2.7,-0.6) .. (2.2,-0.6) -- (0,-0.6) -- cycle;

                \draw[deformed]
                (0,0.6) .. controls (1.1,0.6) and (1.6,0.56) .. (2.2,0.52)
                .. controls (2.7,0.52) and (2.9,0.0) .. (3.25,0.0)
                .. controls (3.6,0.0) and (3.8,0.3) .. (4.3,0.3)
                .. controls (5.0,0.1) and (5.5,-0.1) .. (5.8,-0.2) --
                (5.8,-1.4)
                .. controls (5.5,-1.3) and (5.0,-1.1) .. (4.3,-0.9)
                .. controls (3.8,-0.9) and (3.6,-0.5) .. (3.25,-0.5)
                .. controls (2.9,-0.5) and (2.7,-0.68) .. (2.2,-0.68)
                .. controls (1.6,-0.68) and (1.1,-0.6) .. (0,-0.6) -- cycle;
                \node[primaryblue] at (1.2, 0.0) {$\Omega(\boldsymbol{\mu})$};

                \node[right, red!70!black] at (5.8, -0.6) {$\Gamma_N(\boldsymbol{\mu})$};
                \draw[force] (5.8, -0.8) -- (5.8, -1.4) node[right] {$\vect{h}_N$};
                \draw[->, line width=1.2pt, orange!80!black] (3.6, -0.25) -- (3.6, -0.85) node[right] {$\vect{f}_b$};
                \draw[->, line width=1.0pt, purple!80!black] (5.8, -0.8) -- (6.3, -0.8) node[right] {$\vect{n}$};

                \coordinate (X_ref) at (3.5, 0.0);
                \coordinate (x) at (3.8, -0.3);
                \fill[gray!60] (X_ref) circle (1.2pt);
                \fill[primaryblue] (x) circle (1.5pt) node[right] {$\vect{x} + \vect{u}$};
                \draw[displ] (X_ref) -- (x) node[midway, left] {$\vect{u}$};

                \node[below, font=\bfseries] at (3.0, -1.8) {(b) Deformed Configuration};
            \end{scope}
        \end{tikzpicture}%
    }
    \caption{Kinematic mapping and boundary conditions of the parameterized notched cantilever beam.}
    \label{fig:continuum_kinematics}
\end{figure}

Comparing the undeformed reference geometry $\Omega(\boldsymbol{\mu})$ to the deformed state under applied loads $\vect{f}_b$ and $\vect{h}_N$ (Figure \ref{fig:continuum_kinematics}) demonstrates the displacement field vector $\vect{u}(t)$.

\begin{table}[htbp]
    \centering
    \caption{Definition of kinematic notation and continuum domain boundaries.}
    \label{tab:kinematic_notation}
    \begin{tabular}{l p{10.5cm}}
        \toprule
        \textbf{Symbol}                           & \textbf{Physical Meaning / Definition}                                                                                  \\
        \midrule
        $\Omega(\boldsymbol{\mu})$                & Parameter-dependent physical domain configuration                                                                       \\
        $\Gamma_D(\boldsymbol{\mu})$              & Fixed Dirichlet boundary at left clamped support ($\vect{u} = \vect{0}$)                                                \\
        $\Gamma_N(\boldsymbol{\mu})$              & Neumann traction boundary subjected to dynamic loading                                                                  \\
        $\vect{u}(\vect{x}, t; \boldsymbol{\mu})$ & Spatial displacement field vector at position $\vect{x}$ and time $t$                                                   \\
        $\boldsymbol{\sigma}(\vect{u})$           & Cauchy stress tensor under linear elasticity                                                                            \\
        $\boldsymbol{\varepsilon}(\vect{u})$      & Small-strain tensor: $\boldsymbol{\varepsilon}(\vect{u}) = \frac{1}{2}\left(\nabla\vect{u} + (\nabla\vect{u})^T\right)$ \\
        $\mathbf{D}$                              & Isotropic constitutive elasticity matrix in Voigt notation                                                              \\
        \bottomrule
    \end{tabular}
\end{table}

This formal notation (Table \ref{tab:kinematic_notation}) establishes the continuum framework used throughout the initial configuration and small-strain kinematic relations.

Assuming small strains and linear isotropic elasticity, the strong form governing linear momentum balance on $\Omega(\boldsymbol{\mu}) \times (0, T]$ is:
\begin{equation}
    \underbrace{\rho \ddot{\vect{u}}}_{\text{Inertial Force}} - \underbrace{\nabla \cdot \boldsymbol{\sigma}(\vect{u})}_{\text{Internal Stress Divergence}} = \underbrace{\vect{f}_b}_{\text{Body Force}}
\end{equation}
subject to boundary conditions $\vect{u} = \vect{0}$ on $\Gamma_D(\boldsymbol{\mu})$, traction $\boldsymbol{\sigma}\cdot\vect{n} = \vect{h}_N$ on $\Gamma_N(\boldsymbol{\mu})$, and initial conditions $\vect{u}(\vect{x},0) = \vect{u}_0(\vect{x})$, $\dot{\vect{u}}(\vect{x},0) = \vect{v}_0(\vect{x})$.

% ------------------------------------------------------------------------------
\section{Weak Variational Formulation}
% ------------------------------------------------------------------------------

This section derives the weak variational form of elastodynamics, establishing the kinetic energy, strain energy, and external work forms required for numerical discretization.

Multiplying by kinematically admissible test functions $\delta\vect{u} \in \mathcal{V}_0(\Omega(\boldsymbol{\mu})) = \{ \vect{v} \in [H^1(\Omega(\boldsymbol{\mu}))]^2 \mid \vect{v} = \vect{0} \text{ on } \Gamma_D(\boldsymbol{\mu}) \}$ and applying integration by parts yields the weak form:

\begin{equation}
    \underbrace{m(\ddot{\vect{u}}, \delta\vect{u}; \boldsymbol{\mu})}_{\text{Kinetic Energy Variation}} + \underbrace{k(\vect{u}, \delta\vect{u}; \boldsymbol{\mu})}_{\text{Internal Strain Energy}} = \underbrace{f(\delta\vect{u})}_{\text{External Load Work}} \quad \forall \delta\vect{u} \in \mathcal{V}_0(\Omega(\boldsymbol{\mu}))
\end{equation}

where the bilinear forms and load linear form are defined as:
\begin{align}
    m(\ddot{\vect{u}}, \delta\vect{u}; \boldsymbol{\mu}) & = \int_{\Omega(\boldsymbol{\mu})} \rho \ddot{\vect{u}} \cdot \delta\vect{u} \, d\Omega                                                                      \\
    k(\vect{u}, \delta\vect{u}; \boldsymbol{\mu})        & = \int_{\Omega(\boldsymbol{\mu})} \boldsymbol{\varepsilon}(\delta\vect{u})^T \mat{D} \boldsymbol{\varepsilon}(\vect{u}) \, d\Omega                          \\
    f(\delta\vect{u})                                    & = \int_{\Omega(\boldsymbol{\mu})} \vect{f}_b \cdot \delta\vect{u} \, d\Omega + \int_{\Gamma_N(\boldsymbol{\mu})} \vect{h}_N \cdot \delta\vect{u} \, d\Gamma
\end{align}

% ------------------------------------------------------------------------------
\section{The Shape Parameterization Challenge}
% ------------------------------------------------------------------------------

This section highlights why physical domain variations break conventional snapshot POD extraction, motivating the transition to reference domain pullback integration.

Evaluating $k(\vect{u}, \delta\vect{u}; \boldsymbol{\mu})$ on a varying domain $\Omega(\boldsymbol{\mu})$ introduces two critical computational hurdles:

\begin{table}[htbp]
    \centering
    \caption{Trade-offs between physical remeshing and reference domain pullback mapping.}
    \label{tab:remesh_vs_pullback}
    \resizebox{\textwidth}{!}{%
        \begin{tabular}{p{3.5cm} p{5.5cm} p{5.5cm}}
            \toprule
            \textbf{Criteria}        & \textbf{Naive Physical Remeshing}                           & \textbf{Reference Pullback Mapping}                                \\
            \midrule
            \textbf{Vector Space}    & Inconsistent grid nodes across parameters (Destroys SVD).   & Invariant reference space $\hat{\Omega}$ (Preserves snapshot SVD). \\
            \midrule
            \textbf{Online Assembly} & Full physical remesh \& assembly $\mathcal{O}(E)$ per step. & Precomputed reference integration; scales with active set $E^*$.   \\
            \midrule
            \textbf{Geometry Error}  & Tessellation approximation error.                           & Exact CAD NURBS representation.                                    \\
            \bottomrule
        \end{tabular}%
    }
\end{table}

Evaluating these computational trade-offs (Table \ref{tab:remesh_vs_pullback}) highlights why reference domain pullback mapping is essential for maintaining snapshot SVD consistency.

To guarantee mathematical consistency and achieve sub-millisecond solves, this thesis selects the \textbf{Diffeomorphic Pullback Mapping} approach, shifting parameter dependencies entirely into integration Jacobians on a fixed reference space $\hat{\Omega}$ (detailed in Chapter \ref{ch:pullback}).

% ------------------------------------------------------------------------------
\section{Chapter Summary}
% ------------------------------------------------------------------------------

This section synthesizes the continuum elastodynamic foundation established in this chapter and previews the spatial discretization techniques introduced in Chapter \ref{ch:discretization}.

\begin{tcolorbox}[colback=lightbg, colframe=accentcyan, title=\textbf{Key Takeaway}, fonttitle=\bfseries]
    We established the elastodynamic continuum formulation for the notched cantilever beam and demonstrated that direct physical evaluation on $\Omega(\boldsymbol{\mu})$ breaks state vector compatibility across parameters. In Chapter \ref{ch:discretization}, we discretize this formulation using Isogeometric Analysis (IGA) prior to pulling back integration terms to a reference domain in Chapter \ref{ch:pullback}.
\end{tcolorbox}
